\section{Context}
\label{chap:introduction:context}

The advent of \acs{5G} and \ac{B5G} networks represents an important milestone in telecommunications, introducing new wireless technologies and services that exploit the capabilities of these networks to deliver new functionalities and accomodate growing user demand~\cite{electronics12102200}. These technologies are designed to deliver (i) enhanced capacity, (ii) reduced latency, and (iii) unique levels of programmability, all of which essential to support the needs of different industry verticals~\cite{SensorsPaper}: specific sectors/domains with specific service requirements that are satisfied via tailored \acs{5G}/\acs{B5G} solutions, such as healthcare, automotive and transports~\cite{disruptive5G}. Such services can be offered through Network Applications, software solutions that interact directly with the control planes of \acs{5G}/\acs{B5G} to enable dynamic and network-aware functionalities. This opens the possiblity for \acfp{MNO} to capitalize on new revenue streams by addressing the demands of different verticals~\cite{SEAL5GVerticals}.

Nevertheless, for verticals to fully capitalize on these opportunities, they need to overcome the challenge of integrating their unique requirements with the networks' control and management capabilities, leveraging the programmability and flexibility offered by \acs{5G} networks. Therefore, \acs{3GPP} is heavily invested on the definition of standardized interfaces~\cite{INTR_ETSI_TS_126_531_V18_2_0}. In this context, the \acf{NEF} plays a crucial role, providing secure exposure of network capabilities to third parties, allowing them to interact with the network through standardized interfaces.

However, the \acs{NEF} does not fully address the challenge of integrating network applications with the network's control and management functions, often proving too complex to interact with for vertical developers with no telco expertise. In result, several key standardization initiatives have emerged to address such challenges, including the CAMARA Project, \acs{GSMA} Open Gateway, \acs{CAPIF} and \acs{TMF} to promote interoperability, security, and controlled exposure of network capabilities and resources.

Despite the promising potential of these standardization initiatives and exposure mechanisms, their practical deployment and validation in real-world \acs{5G} environments remains a significant challenge. Consequently, establishing a testbed environment that supports these network exposure interfaces is essential not only to validate their technical feasibility but also to provide a practical platform where vertical developers can experiment with and leverage these capabilities. This dissertation addresses this limitation by implementing, integrating and validating these standardization initiatives and exposure interfaces within the 5GAIner testbed at \acf{ITAv}, culminating with the developement of a Network Application that demonstrates their practical application.


